\chapter{Introducción}

En la actualidad, la contratación de servicios de flete para el traslado de cargas menores, muebles u objetos personales suele realizarse de manera informal, a través de recomendaciones, contactos directos o publicaciones en redes sociales. Esta situación genera diversas problemáticas para los usuarios, tales como la falta de transparencia en los precios, la dificultad para comparar alternativas, la desconfianza respecto al estado de la carga transportada y la ausencia de mecanismos formales para realizar reclamos o verificar la responsabilidad del transportista.

Frente a esta problemática, el presente proyecto propone el diseño de una plataforma tecnológica orientada a conectar clientes que necesitan realizar traslados con transportistas que ofrecen servicios de flete. El sistema busca facilitar la cotización, contratación, seguimiento y evaluación del servicio, incorporando herramientas que permitan mejorar la seguridad, trazabilidad y confianza entre las partes involucradas. Además, se considera un modelo de negocio basado en una comisión porcentual aplicada a cada viaje completado mediante la plataforma.

Los usuarios principales del sistema son tres: el Cliente, el Transportista y el Administrador. El Cliente podrá solicitar servicios de traslado, comparar transportistas disponibles, revisar precios, realizar pagos, seguir el viaje en tiempo real y calificar el servicio recibido. El Transportista podrá registrarse en la plataforma, configurar sus tarifas, recibir solicitudes cercanas, subir evidencia fotográfica del estado de la carga y gestionar sus servicios realizados. Por su parte, el Administrador tendrá la responsabilidad de validar documentación, gestionar usuarios, revisar reportes y supervisar el funcionamiento general de la plataforma.

La naturaleza del software propuesto corresponde a una aplicación móvil para clientes y transportistas, complementada con un panel web administrativo. Esta combinación permite responder a las necesidades operativas de los usuarios que participan directamente en los traslados, mientras se mantiene una gestión centralizada del sistema por parte de la administración. El software se inspira en plataformas colaborativas de intermediación de servicios, adaptadas específicamente al rubro de los fletes y traslados de carga menor.
En este informe se desarrollan los principales aspectos asociados al proyecto de software, comenzando por el análisis de requerimientos del sistema, donde se identifican las necesidades de los usuarios, los requerimientos funcionales y los requerimientos no funcionales. Posteriormente, se presenta el modelado del sistema mediante diagramas BPMN y casos de uso, junto con la justificación del modelo de desarrollo seleccionado. Además, se describen las herramientas CASE utilizadas, las estrategias de soporte, los aspectos de seguridad y calidad del software, y el diseño de mockups que representan la interfaz propuesta para la plataforma.

Finalmente, este proyecto permite aplicar conceptos fundamentales de la Ingeniería de Software, tales como el levantamiento de requerimientos, el modelado de procesos, la planificación del desarrollo, la documentación técnica y el diseño de soluciones centradas en el usuario. De esta manera, la propuesta busca entregar una solución tecnológica viable, organizada y alineada con las necesidades reales de los usuarios involucrados en la contratación y prestación de servicios de flete.

